\begingroup
\let\cleardoublepage\relax


\chapter*{Java Datentypen}
\label{ch:java-datentypen}
\addcontentsline{toc}{chapter}{Java Datentypen}
\endgroup


\section*{Primitive Datentypen}
\addcontentsline{toc}{section}{Primitive Datentypen}

Java arbeitet mit verschiedenen Arten von Daten. Die einfachsten davon heißen "primitive Datentypen".
Diese speichern direkte Werte wie Zahlen oder Ja/Nein-Antworten.

% ============================================================
\newpage


\subsection*{int - Ganzzahlen}
\addcontentsline{toc}{subsection}{int - Ganzzahlen}

\subsubsection*{Was ist int?}\\
Der Datentyp \texttt{int} speichert ganze Zahlen - also Zahlen ohne Komma (wie 5, -10, oder 1000).

\subsubsection*{Wertebereich:}
\begin{itemize}
	\item Von: -2.147.483.648
	\item Bis: 2.147.483.647
	\item Speichergröße: 32 Bit (4 Bytes)
\end{itemize}

\subsubsection*{Syntax:}
\begin{lstlisting}
int variablenName = wert;
\end{lstlisting}

\subsubsection*{Beispiele:}
\begin{lstlisting}
// Deklaration und Initialisierung
int alter = 25;
int temperatur = -5;
int punkte = 0;

// Ohne Initialisierung
int zahl;
zahl = 42;

// Berechnungen
int summe = 10 + 20;          // summe = 30
int differenz = 50 - 15;      // differenz = 35
int produkt = 6 * 7;          // produkt = 42
int quotient = 20 / 4;        // quotient = 5
int rest = 17 % 5;            // rest = 2

// Mit Variablen rechnen
int a = 10;
int b = 3;
int ergebnis = a + b;         // ergebnis = 13

// Inkrement und Dekrement
int zaehler = 5;
zaehler++;                    // zaehler = 6
zaehler--;                    // zaehler = 5

// Zusammengesetzte Operatoren
int wert = 10;
wert += 5;                    // wert = 15 (entspricht: wert = wert + 5)
wert -= 3;                    // wert = 12
wert *= 2;                    // wert = 24
wert /= 4;                    // wert = 6
\end{lstlisting}

\begin{achtung}
	\begin{itemize}
		\item Wenn du eine int-Zahl durch eine int-Zahl teilst, bekommst du wieder eine int-Zahl: \texttt{5 / 2 = 2} (nicht 2.5! Die Nachkommastellen werden einfach abgeschnitten)
		\item Wenn die Zahl zu groß oder zu klein wird, springt sie auf die andere Seite um (das nennt man Überlauf)
		\item In einem Array ist der Anfangswert immer \texttt{0}
		\item Mit int kannst du keine Kommazahlen speichern
	\end{itemize}
\end{achtung}

% ============================================================
\newpage


\subsection*{double - Fließkommazahlen}
\addcontentsline{toc}{subsection}{double - Fließkommazahlen}

\subsubsection*{Was ist double?}\\
Der Datentyp \texttt{double} speichert Kommazahlen - also Zahlen mit Nachkommastellen (wie 3.14, -5.5, oder 0.001).

\subsubsection*{Wertebereich:}
\begin{itemize}
	\item Sehr große und sehr kleine Zahlen mit ca. 15 Dezimalstellen Genauigkeit
	\item Von: $\pm$4.9 $\times$ 10$^{-324}$
	\item Bis: $\pm$1.7 $\times$ 10$^{308}$
	\item Speichergröße: 64 Bit (8 Bytes)
\end{itemize}

\subsubsection*{Syntax:}
\begin{lstlisting}
double variablenName = wert;
\end{lstlisting}

\subsubsection*{Beispiele:}
\begin{lstlisting}
// Deklaration und Initialisierung
double pi = 3.14159;
double temperatur = -5.5;
double gewicht = 75.8;
double ergebnis = 0.0;

// Berechnungen
double summe = 10.5 + 20.3;        // summe = 30.8
double produkt = 2.5 * 4.0;        // produkt = 10.0
double quotient = 7.0 / 2.0;       // quotient = 3.5

// Wichtig: Division mit Nachkommastellen
double a = 5.0;
double b = 2.0;
double teil = a / b;               // teil = 2.5

// int zu double konvertieren
int ganzzahl = 5;
double kommazahl = ganzzahl / 2.0; // kommazahl = 2.5
// ODER
double kommazahl2 = (double) ganzzahl / 2; // kommazahl2 = 2.5

// Wissenschaftliche Notation
double klein = 1.5e-4;             // = 0.00015
double gross = 3.2e8;              // = 320000000.0

// Mathematische Funktionen
double wurzel = Math.sqrt(16.0);   // wurzel = 4.0
double potenz = Math.pow(2, 3);    // potenz = 8.0
double absolut = Math.abs(-5.5);   // absolut = 5.5
\end{lstlisting}

\begin{achtung}
	\begin{itemize}
		\item In Java schreibst du einen Punkt statt Komma: \texttt{3.14} ist richtig, \texttt{3,14} ist falsch
		\item Computer können nicht immer exakt rechnen: manchmal kommt bei \texttt{0.1 + 0.2} nicht genau \texttt{0.3} heraus
		\item Wichtig: \texttt{5 / 2} ergibt \texttt{2}, aber \texttt{5.0 / 2.0} ergibt \texttt{2.5} - mindestens eine Zahl muss eine Kommazahl sein
		\item In einem Array ist der Anfangswert immer \texttt{0.0}
		\item Für Geld solltest du andere Datentypen verwenden (z.B. \texttt{BigDecimal}), weil Rundungsfehler problematisch sein können
	\end{itemize}
\end{achtung}

% ============================================================
\newpage


\subsection*{boolean - Wahrheitswerte}
\addcontentsline{toc}{subsection}{boolean - Wahrheitswerte}

\subsubsection*{Was ist boolean?}\\
Der Datentyp \texttt{boolean} speichert Wahrheitswerte: entweder \texttt{true} (wahr) oder \texttt{false} (falsch).

\subsubsection*{Wertebereich:}
\begin{itemize}
	\item Nur zwei mögliche Werte: \texttt{true} oder \texttt{false}
	\item Speichergröße: 1 Bit (theoretisch), praktisch abhängig von JVM
\end{itemize}

\subsubsection*{Syntax:}
\begin{lstlisting}
boolean variablenName = true/false;
\end{lstlisting}

\subsubsection*{Beispiele:}
\begin{lstlisting}
// Deklaration und Initialisierung
boolean istWahr = true;
boolean istFalsch = false;
boolean hatBestanden = true;

// Vergleichsoperatoren ergeben boolean
int alter = 18;
boolean istErwachsen = (alter >= 18);      // true
boolean istKind = (alter < 12);            // false
boolean istGenau18 = (alter == 18);        // true
boolean istNicht18 = (alter != 18);        // false

// Logische Operatoren
boolean a = true;
boolean b = false;

boolean und = a && b;                      // false (beide muessen true sein)
boolean oder = a || b;                     // true (mindestens einer true)
boolean nicht = !a;                        // false (Negation)

// In Bedingungen
int punkte = 85;
boolean bestanden = punkte >= 50;
if (bestanden) {
    System.out.println("Bestanden!");
}

// Kombinierte Bedingungen
int temperatur = 25;
boolean angenehm = temperatur >= 18 && temperatur <= 28;

// Boolean in Schleifen
boolean weitermachen = true;
while (weitermachen) {
    // Code
    weitermachen = false;  // Schleife beenden
}

// Vergleich mit Strings
String name = "Max";
boolean istMax = name.equals("Max");       // true
\end{lstlisting}

\begin{achtung}
	\begin{itemize}
		\item Es gibt nur zwei Werte: \texttt{true} (wahr) oder \texttt{false} (falsch) - keine Zahlen!
		\item Alles klein schreiben: \texttt{true} und \texttt{false}, nicht \texttt{True} oder \texttt{False}
		\item In einem Array ist der Anfangswert immer \texttt{false}
		\item Zum Vergleichen: Bei Zahlen \texttt{==}, bei Text \texttt{.equals()}
	\end{itemize}
\end{achtung}

% ============================================================
\newpage


\subsection*{char - Einzelnes Zeichen}
\addcontentsline{toc}{subsection}{char - Einzelnes Zeichen}

\subsubsection*{Was ist char?}\\
Der Datentyp \texttt{char} speichert genau ein einzelnes Zeichen (wie 'A', '5' oder '!').

\subsubsection*{Wertebereich:}
\begin{itemize}
	\item Alle Unicode-Zeichen (Buchstaben, Ziffern, Sonderzeichen)
	\item Von: 0 bis 65.535 (als Zahlenwert)
	\item Speichergröße: 16 Bit (2 Bytes)
\end{itemize}

\subsubsection*{Syntax:}
\begin{lstlisting}
char variablenName = 'Zeichen';  // EINFACHE Anfuehrungszeichen!
\end{lstlisting}

\subsubsection*{Beispiele:}
\begin{lstlisting}
// Deklaration und Initialisierung
char buchstabe = 'A';
char ziffer = '5';
char sonderzeichen = '!';
char leerzeichen = ' ';

// Mit Unicode-Werten
char euro = '\u20AC';                      // €-Symbol
char herzchen = '\u2764';                  // ❤

// Escape-Sequenzen
char neuzeile = '\n';                      // Zeilenumbruch
char tab = '\t';                           // Tabulator
char backslash = '\\';                     // \
char anfuehrung = '\'';                    // '

// Zeichen vergleichen
char note = 'A';
if (note == 'A') {
    System.out.println("Sehr gut!");
}

// Zeichen in Berechnungen (ASCII/Unicode-Werte)
char a = 'A';                              // Unicode-Wert: 65
char b = 'B';                              // Unicode-Wert: 66
int differenz = b - a;                     // differenz = 1

// Kleinbuchstabe zu Großbuchstabe
char klein = 'a';                          // Unicode-Wert: 97
char gross = (char) (klein - 32);          // gross = 'A' (Wert: 65)

// Prüfen ob Buchstabe
char zeichen = 'M';
boolean istBuchstabe = (zeichen >= 'A' && zeichen <= 'Z') ||
                       (zeichen >= 'a' && zeichen <= 'z');

// Prüfen ob Ziffer
char z = '7';
boolean istZiffer = (z >= '0' && z <= '9');

// char aus String extrahieren
String text = "Hallo";
char erstes = text.charAt(0);              // erstes = 'H'
\end{lstlisting}

\begin{achtung}
	\begin{itemize}
		\item Immer einfache Anführungszeichen verwenden: \texttt{'A'} ist richtig, \texttt{"A"} ist falsch
		\item Nur genau ein Zeichen erlaubt: \texttt{'A'} funktioniert, \texttt{'AB'} nicht
		\item \texttt{char} ist etwas anderes als \texttt{String} - ein char ist nur ein einzelnes Zeichen
		\item In einem Array ist der Anfangswert das "Null-Zeichen" (ein unsichtbares Zeichen)
		\item Jedes Zeichen hat intern eine Nummer - deshalb kannst du mit char auch rechnen
	\end{itemize}
\end{achtung}

% ============================================================
\newpage


\subsection*{byte - Sehr kleine Ganzzahlen}
\addcontentsline{toc}{subsection}{byte - Sehr kleine Ganzzahlen}

\subsubsection*{Was ist byte?}\\
Der Datentyp \texttt{byte} speichert sehr kleine ganze Zahlen.
Wird selten benutzt, meist nur wenn man Speicherplatz sparen muss.

\subsubsection*{Wertebereich:}
\begin{itemize}
	\item Von: -128
	\item Bis: 127
	\item Speichergröße: 8 Bit (1 Byte)
\end{itemize}

\subsubsection*{Syntax:}
\begin{lstlisting}
byte variablenName = wert;
\end{lstlisting}

\subsubsection*{Beispiele:}
\begin{lstlisting}
// Deklaration und Initialisierung
byte alter = 25;
byte temperatur = -15;
byte prozent = 100;

// Häufig bei Datei-Operationen
byte[] daten = new byte[1024];

// Konvertierung von int zu byte
int gross = 100;
byte klein = (byte) gross;

// Vorsicht bei Überschreitung
byte max = 127;
byte ueberlauf = (byte) (max + 1);         // ueberlauf = -128 (Overflow!)
\end{lstlisting}

\begin{achtung}
	\begin{itemize}
		\item Sehr kleiner Wertebereich: nur -128 bis 127
		\item Wenn die Zahl zu groß wird, springt sie um (Überlauf)
		\item Wird meistens nur für spezielle Dinge verwendet
		\item In einem Array ist der Anfangswert immer \texttt{0}
	\end{itemize}
\end{achtung}

% ============================================================
\newpage


\subsection*{short - Kleine Ganzzahlen}
\addcontentsline{toc}{subsection}{short - Kleine Ganzzahlen}

\subsubsection*{Was ist short?}\\
Der Datentyp \texttt{short} speichert kleine ganze Zahlen.
Ein Mittelweg zwischen \texttt{byte} und \texttt{int}.

\subsubsection*{Wertebereich:}
\begin{itemize}
	\item Von: -32.768
	\item Bis: 32.767
	\item Speichergröße: 16 Bit (2 Bytes)
\end{itemize}

\subsubsection*{Syntax:}
\begin{lstlisting}
short variablenName = wert;
\end{lstlisting}

\subsubsection*{Beispiele:}
\begin{lstlisting}
// Deklaration und Initialisierung
short jahr = 2024;
short distanz = 5000;

// Konvertierung
int gross = 1000;
short klein = (short) gross;
\end{lstlisting}

\begin{achtung}
	\begin{itemize}
		\item Wird selten verwendet, meist wird direkt \texttt{int} genommen
		\item Wertebereich: -32.768 bis 32.767
		\item Standard-Wert bei Arrays: \texttt{0}
	\end{itemize}
\end{achtung}

% ============================================================
\newpage


\subsection*{long - Sehr große Ganzzahlen}
\addcontentsline{toc}{subsection}{long - Sehr große Ganzzahlen}

\subsubsection*{Was ist long?}\\
Der Datentyp \texttt{long} speichert sehr große ganze Zahlen - noch größer als int.

\subsubsection*{Wertebereich:}
\begin{itemize}
	\item Von: -9.223.372.036.854.775.808
	\item Bis: 9.223.372.036.854.775.807
	\item Speichergröße: 64 Bit (8 Bytes)
\end{itemize}

\subsubsection*{Syntax:}
\begin{lstlisting}
long variablenName = wert;
long grosseZahl = wertL;  // L am Ende fuer long-Literal
\end{lstlisting}

\subsubsection*{Beispiele:}
\begin{lstlisting}
// Deklaration und Initialisierung
long weltbevoelkerung = 8000000000L;       // L am Ende!
long millisekunden = System.currentTimeMillis();

// Ohne L wird es als int interpretiert (Fehler bei großen Zahlen)
// long falsch = 8000000000;               // FEHLER!
long richtig = 8000000000L;                // Korrekt

// Berechnungen
long summe = 1000000000L + 2000000000L;
\end{lstlisting}

\begin{achtung}
	\begin{itemize}
		\item Bei großen Zahlen musst du ein \texttt{L} ans Ende schreiben: \texttt{1000000000L}
		\item Ohne \texttt{L} denkt Java, es ist eine int-Zahl, und das kann schiefgehen
		\item In einem Array ist der Anfangswert immer \texttt{0L}
		\item Benutze long, wenn deine Zahlen größer als etwa 2 Milliarden werden können
	\end{itemize}
\end{achtung}

% ============================================================
\newpage


\subsection*{float - Einfache Fließkommazahlen}
\addcontentsline{toc}{subsection}{float - Einfache Fließkommazahlen}

\subsubsection*{Was ist float?}\\
Der Datentyp \texttt{float} speichert Fließkommazahlen mit einfacher Genauigkeit.
Weniger genau als \texttt{double}.

\subsubsection*{Wertebereich:}
\begin{itemize}
	\item Ca. 6-7 Dezimalstellen Genauigkeit
	\item Von: $\pm$1.4 $\times$ 10$^{-45}$
	\item Bis: $\pm$3.4 $\times$ 10$^{38}$
	\item Speichergröße: 32 Bit (4 Bytes)
\end{itemize}

\subsubsection*{Syntax:}
\begin{lstlisting}
float variablenName = wertF;  // F am Ende fuer float-Literal
\end{lstlisting}

\subsubsection*{Beispiele:}
\begin{lstlisting}
// Deklaration und Initialisierung
float pi = 3.14F;                          // F am Ende!
float temperatur = 36.5F;
float gewicht = 75.8F;

// Ohne F wird es als double interpretiert
// float falsch = 3.14;                    // Warnung/Fehler
float richtig = 3.14F;                     // Korrekt

// Berechnungen
float summe = 10.5F + 20.3F;
\end{lstlisting}

\begin{achtung}
	\begin{itemize}
		\item \texttt{F} oder \texttt{f} am Ende anhängen: \texttt{3.14F}
		\item Weniger genau als \texttt{double}
		\item In der Praxis wird meist \texttt{double} bevorzugt
		\item Standard-Wert bei Arrays: \texttt{0.0F}
	\end{itemize}
\end{achtung}

% ============================================================
\newpage


\section*{String - Zeichenketten}
\addcontentsline{toc}{section}{String - Zeichenketten}

\subsection*{Was ist ein String?}

\subsubsection*{Beschreibung}\\
Ein \texttt{String} speichert Text - also mehrere Zeichen zusammen (wie "Hallo" oder "Max Mustermann").
String ist etwas Besonderes: Er gehört nicht zu den primitiven Datentypen.

\subsubsection*{Eigenschaften:}
\begin{itemize}
	\item Wenn du einen String einmal erstellt hast, kannst du ihn nicht mehr ändern - nur einen neuen erstellen
	\item Strings schreibst du mit doppelten Anführungszeichen: \texttt{"Text"}
	\item Ein String kann so lang sein wie du willst
\end{itemize}

\subsubsection*{Syntax:}
\begin{lstlisting}
String variablenName = "Text";
\end{lstlisting}

\subsubsection*{Beispiele:}
\begin{lstlisting}
// Deklaration und Initialisierung
String name = "Max Mustermann";
String begruessung = "Hallo Welt!";
String leer = "";                          // Leerer String
String mitZahl = "Ich bin 25 Jahre alt";

// String-Konkatenation (Verkettung)
String vorname = "Max";
String nachname = "Mustermann";
String vollername = vorname + " " + nachname;  // "Max Mustermann"

// Mit Zahlen verketten
int alter = 25;
String text = "Ich bin " + alter + " Jahre alt";

// String-Methoden
String wort = "Hallo";
int laenge = wort.length();                // 5
char erstes = wort.charAt(0);              // 'H'
String gross = wort.toUpperCase();         // "HALLO"
String klein = wort.toLowerCase();         // "hallo"

// Strings vergleichen
String a = "Hallo";
String b = "Hallo";
boolean gleich = a.equals(b);              // true
boolean gleichIgnoreCase = a.equalsIgnoreCase("HALLO");  // true

// WICHTIG: Nicht mit == vergleichen!
String s1 = "Test";
String s2 = new String("Test");
boolean falsch = (s1 == s2);               // false (vergleicht Referenzen)
boolean richtig = s1.equals(s2);           // true (vergleicht Inhalt)

// Teilstrings
String text2 = "Hallo Welt";
String teil = text2.substring(0, 5);       // "Hallo"
String ab = text2.substring(6);            // "Welt"

// Suchen in Strings
String satz = "Java ist toll";
boolean enthaelt = satz.contains("ist");   // true
int position = satz.indexOf("ist");        // 5
boolean beginnt = satz.startsWith("Java"); // true
boolean endet = satz.endsWith("toll");     // true

// Ersetzen
String original = "Hallo Welt";
String neu = original.replace("Welt", "Java");  // "Hallo Java"

// Leerzeichen entfernen
String mitLeer = "  Text  ";
String ohneLeer = mitLeer.trim();          // "Text"

// String in Array umwandeln
String liste = "Apfel,Birne,Orange";
String[] fruechte = liste.split(",");      // ["Apfel", "Birne", "Orange"]

// Zahlen zu String
int zahl = 42;
String zahlString = String.valueOf(zahl);  // "42"
// ODER
String zahlString2 = "" + zahl;            // "42"

// String zu Zahl
String nummerText = "123";
int nummer = Integer.parseInt(nummerText); // 123
double komma = Double.parseDouble("3.14"); // 3.14
\end{lstlisting}

\begin{achtung}
	\begin{itemize}
		\item Doppelte Anführungszeichen verwenden: \texttt{"Text"} ist richtig, \texttt{'Text'} ist falsch
		\item Zum Vergleichen von Strings immer \texttt{.equals()} benutzen, nicht \texttt{==}
		\item String schreibst du immer mit großem S am Anfang
		\item Wenn du einen String ändern willst, bekommst du eigentlich einen neuen String zurück
		\item \texttt{null} bedeutet "gar kein Text", \texttt{""} ist "leerer Text" - das ist ein Unterschied
		\item In einem Array ist der Anfangswert immer \texttt{null}
	\end{itemize}
\end{achtung}

% ============================================================
\newpage


\section*{Typkonvertierung}
\addcontentsline{toc}{section}{Typkonvertierung}

\subsection*{Implizite Konvertierung (Automatisch)}

\subsubsection*{Beschreibung}\\
Java wandelt automatisch kleinere Zahlen in größere um. Zum Beispiel wird eine int-Zahl automatisch zu double, wenn nötig.

\subsubsection*{Beispiele:}
\begin{lstlisting}
// int zu double (automatisch)
int ganzzahl = 5;
double kommazahl = ganzzahl;               // kommazahl = 5.0

// byte zu int (automatisch)
byte klein = 100;
int gross = klein;                         // gross = 100

// char zu int (automatisch)
char buchstabe = 'A';
int code = buchstabe;                      // code = 65 (Unicode-Wert)

// In Berechnungen
int a = 5;
double b = 2.5;
double ergebnis = a + b;                   // ergebnis = 7.5
\end{lstlisting}

\subsection*{Explizite Konvertierung (Type Cast)}

\subsubsection*{Beschreibung}\\
Wenn du eine größere Zahl in eine kleinere umwandeln willst, musst du Java das ausdrücklich sagen.
Dabei können Informationen verloren gehen!

\subsubsection*{Syntax:}
\begin{lstlisting}
zielTyp variable = (zielTyp) quellVariable;
\end{lstlisting}

\subsubsection*{Beispiele:}
\begin{lstlisting}
// double zu int (Nachkommastellen gehen verloren!)
double kommazahl = 3.7;
int ganzzahl = (int) kommazahl;            // ganzzahl = 3 (nicht 4!)

// long zu int
long gross = 1000L;
int klein = (int) gross;

// int zu byte (Vorsicht bei großen Zahlen!)
int zahl = 200;
byte kleinByte = (byte) zahl;              // Overflow!

// Division mit Cast
int a = 5;
int b = 2;
double ergebnis = (double) a / b;          // ergebnis = 2.5

// Runden statt Abschneiden
double wert = 3.7;
int gerundet = (int) Math.round(wert);     // gerundet = 4
\end{lstlisting}

\begin{achtung}
	\begin{itemize}
		\item Wenn du double zu int umwandelst, werden die Nachkommastellen einfach weggeschnitten, nicht gerundet! (3.9 wird zu 3, nicht zu 4)
		\item Wenn die Zahl zu groß ist, können seltsame Werte herauskommen
		\item Zum Runden benutze \texttt{Math.round()}
		\item Automatisch: klein zu groß (z.B. int zu double)
		\item Mit Klammern: groß zu klein (z.B. double zu int)
	\end{itemize}
\end{achtung}

% ============================================================
\newpage


\section*{Übersicht: Primitive Datentypen}
\addcontentsline{toc}{section}{Übersicht: Primitive Datentypen}

\begin{table}[h]
	\centering
	\begin{tabular}{|l|l|l|l|}
		\hline
		\textbf{Typ} & \textbf{Größe} & \textbf{Wertebereich} & \textbf{Standard} \\
		\hline
		byte & 8 Bit & -128 bis 127 & 0 \\
		\hline
		short & 16 Bit & -32.768 bis 32.767 & 0 \\
		\hline
		int & 32 Bit & -2.147.483.648 bis 2.147.483.647 & 0 \\
		\hline
		long & 64 Bit & ca. -9,2 $\times$ 10$^{18}$ bis 9,2 $\times$ 10$^{18}$ & 0L \\
		\hline
		float & 32 Bit & ca. $\pm$3,4 $\times$ 10$^{38}$ & 0.0F \\
		\hline
		double & 64 Bit & ca. $\pm$1,7 $\times$ 10$^{308}$ & 0.0 \\
		\hline
		char & 16 Bit & 0 bis 65.535 (Unicode) & \('\)\textbackslash u0000' \\
		\hline
		boolean & 1 Bit & true oder false & false \\
		\hline
	\end{tabular}
	\caption{Übersicht der primitiven Datentypen in Java}
\end{table}


\section*{Arrays}
\addcontentsline{toc}{section}{Arrays}

\subsection*{Was ist ein Array?}

\subsubsection*{Beschreibung}\\
Ein Array ist wie eine Liste mit nummerierten Fächern. Jedes Fach kann einen Wert speichern.
Wichtig: Alle Werte müssen vom gleichen Typ sein (z.B. nur Zahlen oder nur Text). Die Anzahl der Fächer ist fest.

\subsubsection*{Eigenschaften:}
\begin{itemize}
	\item Die Größe ist fest: Wenn du ein Array mit 5 Plätzen erstellst, bleiben es 5 Plätze
	\item Die Nummerierung beginnt bei 0: Das erste Fach hat die Nummer 0, das letzte die Nummer \texttt{length - 1}
	\item Alle Werte im Array müssen vom gleichen Typ sein
	\item Neue Arrays sind automatisch mit Startwerten gefüllt (0 für Zahlen, false für boolean, null für Text)
\end{itemize}

\subsubsection*{Syntax:}
\begin{lstlisting}
// Deklaration und Initialisierung mit Größe
datentyp[] arrayName = new datentyp[groesse];

// Deklaration und Initialisierung mit Werten
datentyp[] arrayName = {wert1, wert2, wert3};
\end{lstlisting}

\subsubsection*{Beispiele:}
\begin{lstlisting}
// Array mit fester Größe erstellen
int[] zahlen = new int[5];        // 5 Plätze, alle 0
String[] namen = new String[3];   // 3 Plätze, alle null
boolean[] flags = new boolean[4]; // 4 Plätze, alle false

// Array mit Werten erstellen
int[] noten = {1, 2, 3, 2, 1};
String[] wochentage = {"Mo", "Di", "Mi", "Do", "Fr"};
double[] preise = {9.99, 14.50, 7.25};

// Mehrdimensionale Arrays
int[][] matrix = new int[3][3];   // 3x3 Matrix
int[][] tabelle = {{1, 2}, {3, 4}, {5, 6}};
\end{lstlisting}

\begin{achtung}
	\begin{itemize}
		\item Die Größe kannst du nach dem Erstellen nicht mehr ändern
		\item Die Nummerierung beginnt bei 0, nicht bei 1!
		\item Alle Werte müssen vom gleichen Typ sein
		\item Wenn du ein Array noch nicht erstellt hast, ist es \texttt{null}
	\end{itemize}
\end{achtung}

% ============================================================
\newpage


\subsection*{array.length - Länge eines Arrays}
\addcontentsline{toc}{subsection}{array.length - Länge eines Arrays}

\subsubsection*{Was ist array.length?}\\
Mit \texttt{length} kannst du herausfinden, wie viele Plätze das Array hat. Wichtig: Bei Arrays schreibst du \textbf{keine} Klammern dahinter!

\subsubsection*{Eigenschaften:}
\begin{itemize}
	\item Ist ein Attribut, keine Methode: \textbf{ohne} Klammern \texttt{()}
	\item Gibt die Größe zurück, mit der das Array erstellt wurde
	\item Funktioniert auch bei leeren Arrays (nur Standard-Werte)
\end{itemize}

\subsubsection*{Syntax:}
\begin{lstlisting}
int laenge = arrayName.length;  // OHNE Klammern!
\end{lstlisting}

\subsubsection*{Beispiele:}
\begin{lstlisting}
// Länge verschiedener Arrays
int[] zahlen = {10, 20, 30, 40, 50};
int anzahl = zahlen.length;               // anzahl = 5

String[] namen = new String[10];
int groesse = namen.length;               // groesse = 10 (auch wenn leer!)

// Leeres Array
int[] leer = new int[0];
int laenge = leer.length;                 // laenge = 0

// In Schleife verwenden
int[] werte = {1, 2, 3, 4};
for (int i = 0; i < werte.length; i++) {
    System.out.println(werte[i]);
}

// Letztes Element zugreifen
int[] liste = {100, 200, 300};
int letzter = liste[liste.length - 1];   // letzter = 300

// Bei mehrdimensionalen Arrays
int[][] matrix = new int[3][4];
int zeilen = matrix.length;               // zeilen = 3
int spalten = matrix[0].length;           // spalten = 4
\end{lstlisting}

\begin{achtung}
	\begin{itemize}
		\item Bei Arrays: \texttt{.length} ohne Klammern
		\item Bei Strings: \texttt{.length()} mit Klammern - das ist ein häufiger Fehler!
		\item Gibt die Anzahl der Plätze zurück, nicht die Anzahl der gefüllten Plätze
		\item Auch wenn alle Werte noch 0 oder null sind, hat das Array trotzdem seine feste Länge
	\end{itemize}
\end{achtung}

% ============================================================
\newpage


\subsection*{array[index] - Element lesen}
\addcontentsline{toc}{subsection}{array[index] - Element lesen}

\subsubsection*{Was ist array[index]?}\\
Mit eckigen Klammern \texttt{[]} kannst du auf ein bestimmtes Fach im Array zugreifen.
Die Zahl in den Klammern ist die Nummer des Fachs (beginnt bei 0).

\subsubsection*{Eigenschaften:}
\begin{itemize}
	\item Index beginnt bei 0: Erstes Element ist \texttt{array[0]}
	\item Letztes Element ist \texttt{array[array.length - 1]}
	\item Bei ungültigem Index: \texttt{ArrayIndexOutOfBoundsException}
\end{itemize}

\subsubsection*{Syntax:}
\begin{lstlisting}
datentyp element = arrayName[index];
\end{lstlisting}

\subsubsection*{Beispiele:}
\begin{lstlisting}
// Einzelne Elemente lesen
int[] zahlen = {10, 20, 30, 40, 50};
int erstes = zahlen[0];                   // erstes = 10
int zweites = zahlen[1];                  // zweites = 20
int letztes = zahlen[4];                  // letztes = 50

// Alle Elemente durchlaufen
String[] woerter = {"Hallo", "Welt", "Java"};
for (int i = 0; i < woerter.length; i++) {
    String wort = woerter[i];
    System.out.println(wort);
}

// Bestimmte Position prüfen
int[] noten = {1, 2, 3, 2, 1};
if (noten[2] == 3) {
    System.out.println("Dritte Note ist eine 3");
}

// Letztes Element
int[] liste = {5, 10, 15, 20};
int letzterWert = liste[liste.length - 1];  // letzterWert = 20

// Element in Berechnung verwenden
int[] preise = {100, 200, 150};
int summe = preise[0] + preise[1] + preise[2];  // summe = 450

// Mehrdimensionale Arrays
int[][] matrix = {{1, 2, 3}, {4, 5, 6}};
int wert = matrix[0][1];                  // wert = 2 (1. Zeile, 2. Spalte)
\end{lstlisting}

\begin{achtung}
	\begin{itemize}
		\item Die Nummerierung beginnt bei 0, nicht bei 1!
		\item Die Nummer muss zwischen 0 und \texttt{array.length - 1} liegen
		\item Wenn die Nummer zu groß oder zu klein ist, stürzt das Programm ab
		\item Das letzte Fach erreichst du mit \texttt{array[array.length - 1]}
	\end{itemize}
\end{achtung}

% ============================================================
\newpage


\subsection*{array[index] = wert - Element schreiben}
\addcontentsline{toc}{subsection}{array[index] = wert - Element schreiben}

\subsubsection*{Was ist array[index] = wert?}\\
So kannst du einen Wert in ein bestimmtes Fach des Arrays legen (oder den alten Wert überschreiben).

\subsubsection*{Eigenschaften:}
\begin{itemize}
	\item Überschreibt den vorherigen Wert an dieser Position
	\item Der neue Wert muss zum Datentyp des Arrays passen
	\item Index muss gültig sein (0 bis \texttt{length - 1})
\end{itemize}

\subsubsection*{Syntax:}
\begin{lstlisting}
arrayName[index] = neuerWert;
\end{lstlisting}

\subsubsection*{Beispiele:}
\begin{lstlisting}
// Array befüllen
int[] zahlen = new int[5];                // Array mit 5 Plätzen, alle 0
zahlen[0] = 10;                           // Erstes Element auf 10 setzen
zahlen[1] = 20;                           // Zweites Element auf 20 setzen
zahlen[4] = 50;                           // Letztes Element auf 50 setzen
// Ergebnis: {10, 20, 0, 0, 50}

// Werte ändern
String[] namen = {"Anna", "Bob", "Clara"};
namen[1] = "Max";                         // "Bob" wird zu "Max"
// Ergebnis: {"Anna", "Max", "Clara"}

// In Schleife befüllen
int[] quadrate = new int[5];
for (int i = 0; i < quadrate.length; i++) {
    quadrate[i] = i * i;
}
// Ergebnis: {0, 1, 4, 9, 16}

// Wert erhöhen
int[] punkte = {10, 20, 30};
punkte[0] = punkte[0] + 5;                // Erstes Element um 5 erhöhen
// ODER kürzer:
punkte[0] += 5;
// Ergebnis: {15, 20, 30}

// Array "leeren" (auf Standard-Wert setzen)
int[] werte = {1, 2, 3, 4, 5};
for (int i = 0; i < werte.length; i++) {
    werte[i] = 0;
}
// Ergebnis: {0, 0, 0, 0, 0}

// Zwei Elemente tauschen
int[] liste = {10, 20, 30};
int temp = liste[0];                      // Zwischenspeichern
liste[0] = liste[2];                      // 30 nach vorne
liste[2] = temp;                          // 10 nach hinten
// Ergebnis: {30, 20, 10}

// Array aus Benutzereingabe füllen
int[] zahlenListe = new int[3];
for (int i = 0; i < zahlenListe.length; i++) {
    // zahlenListe[i] = scanner.nextInt();
}

// Mehrdimensionale Arrays
int[][] matrix = new int[2][2];
matrix[0][0] = 1;
matrix[0][1] = 2;
matrix[1][0] = 3;
matrix[1][1] = 4;
// Ergebnis: {{1, 2}, {3, 4}}
\end{lstlisting}

\begin{achtung}
	\begin{itemize}
		\item Die Nummer muss zwischen 0 und \texttt{array.length - 1} liegen
		\item Wenn die Nummer ungültig ist, stürzt das Programm ab
		\item Der Wert muss zum Typ des Arrays passen (Zahl für int-Array, Text für String-Array usw.)
		\item Der alte Wert an dieser Stelle geht verloren - er wird komplett ersetzt
	\end{itemize}
\end{achtung}

% ============================================================
\newpage


\section*{Häufige Fehler}
\addcontentsline{toc}{section}{Häufige Fehler}

\subsection*{String mit == vergleichen}

\begin{lstlisting}
// FALSCH
String a = "Hallo";
String b = "Hallo";
if (a == b) { }  // Vergleicht Referenzen, nicht Inhalt!

// RICHTIG
if (a.equals(b)) { }  // Vergleicht Inhalt
\end{lstlisting}

\subsection*{Integer-Division}

\begin{lstlisting}
// FALSCH - ergibt 2, nicht 2.5
int a = 5;
int b = 2;
double ergebnis = a / b;  // ergebnis = 2.0

// RICHTIG
double ergebnis = (double) a / b;  // ergebnis = 2.5
// ODER
double ergebnis = 5.0 / 2.0;  // ergebnis = 2.5
\end{lstlisting}

\subsection*{char vs. String verwechseln}

\begin{lstlisting}
// FALSCH
char buchstabe = "A";  // Fehler: doppelte Anfuehrungszeichen
String wort = 'Hallo';  // Fehler: einfache Anfuehrungszeichen

// RICHTIG
char buchstabe = 'A';   // einfache Anfuehrungszeichen
String wort = "Hallo";  // doppelte Anfuehrungszeichen
\end{lstlisting}

\subsection*{L und F vergessen}

\begin{lstlisting}
// FALSCH
long gross = 5000000000;  // Fehler: zu gross fuer int
float zahl = 3.14;        // Warnung: double wird zu float

// RICHTIG
long gross = 5000000000L;  // L am Ende
float zahl = 3.14F;        // F am Ende
\end{lstlisting}

\subsection*{Overflow nicht beachten}

\begin{lstlisting}
// Problem
byte b = 127;
b++;  // b = -128 (Overflow!)

int i = 2147483647;
i++;  // i = -2147483648 (Overflow!)
\end{lstlisting}

\subsection*{null vs. leerer String}

\begin{lstlisting}
String a = null;      // Keine Referenz (keine Methoden aufrufbar!)
String b = "";        // Leerer String (Methoden aufrufbar)

// a.length();        // FEHLER: NullPointerException
int laenge = b.length();  // Funktioniert: laenge = 0
\end{lstlisting}

\subsection*{String vs. Array - length}

\begin{lstlisting}
// FALSCH
String text = "Hallo";
int laenge = text.length;  // Fehler: Strings brauchen ()

int[] zahlen = {1, 2, 3};
int anzahl = zahlen.length();  // Fehler: Arrays haben KEINE ()

// RICHTIG
int laenge = text.length();    // String MIT Klammern
int anzahl = zahlen.length;    // Array OHNE Klammern
\end{lstlisting}

\subsection*{Array-Index-Grenzen}

\begin{lstlisting}
// FALSCH
int[] zahlen = {10, 20, 30};  // length = 3, Indizes: 0, 1, 2
int x = zahlen[3];            // Fehler! Index 3 existiert nicht
int y = zahlen[-1];           // Fehler! Negative Indizes nicht erlaubt

// RICHTIG
int letzter = zahlen[zahlen.length - 1];  // Index 2
int erster = zahlen[0];                   // Index 0
\end{lstlisting}

\subsection*{Array-Größe nach Erstellung ändern}

\begin{lstlisting}
// FALSCH
int[] zahlen = new int[5];
zahlen.length = 10;  // Fehler! length ist read-only

// RICHTIG - neues Array erstellen und Werte kopieren
int[] zahlen = new int[5];
int[] neuesArray = new int[10];
for (int i = 0; i < zahlen.length; i++) {
    neuesArray[i] = zahlen[i];
}
zahlen = neuesArray;  // Referenz auf neues Array
\end{lstlisting}
